\pagelayout{wide}
\chapter*{Summary and Conclusion\,\,}
\labch{conclusion}
\addcontentsline{toc}{chapter}{Conclusion}

I have demonstrated various different methods of the \ac{SIM} reconstruction from images of an acquisition using
    structured illumination.
Two algorithms for reconstruction, the Wiener filter and the optimization reconstructions in several forms and
    modifications were presented, implemented and their results where presented compared and a commentary was given as 
    to their drawbacks and benefits.
There were given options and prospects for improvement upon these algorithms that should lead to further enhancing the
    results achieved.
Also, a demonstration of the possibility of use of a reduced number of acquisitions was made, that was commented upon in
    the literature, but the results were previously not presented.

A faithful image formation model was developed in the context of the pursuit of estimating the transfer function of the 
    optical system.
The theoretical sources of noise were discussed, before a model of the noise was stated and its validity was indicated.
Then a procedure for estimating the \ac{PSF} from a specific acquisition of micro-spheres (beads) was presented and
    realized.
Then a model of the source of the acquisition was developed via a custom algorithm inspired by other image processing
    algorithms with a result that is well justified by the sensed image.
Two other possible methods of estimating the transfer function were presented and implemented, one inspired by the 
    optimization reconstruction algorithm, with the use of the source model developed and the other also via the source 
    model, but by use of the \ac{LR} image sum.
The averaging of beads estimation was realized on different patches of the image plane and comparison of the
    reconstructions made using the estimates was demonstrated.
Finally the validity and justification for use of these estimations was demonstrated by a reconstruction of the cell
    using the estimate.
A thorough discussion was given on the possible improvements upon the algorithms used.

The final reconstruction of the cell using the estimate demonstrates a compelling result with a visible improvement in
    the quality, artefacts partially eliminated and the structure of the \acp{CCP} on the cell membrane surface
    visible in better contrast against the background noise.
I have demonstrated the validity of the procedures developed and the achieved results improvement is unmistakable.
Future possibilities for and paths towards improvement are clearly stated.
\pagelayout{margin}
